\documentclass[10pt,a4paper]{amsart}
\usepackage[margin=19mm]{geometry}
\usepackage{amsmath,amssymb,mathtools}
\usepackage[scr=rsfs,cal=euler]{mathalfa}
\usepackage{xfrac}
\usepackage[unicode,hidelinks,bookmarks=false]{hyperref}
\newcommand{\ie}{i.e.\ }
\newcommand{\cf}{cf.\ }
\newcommand{\resp}{respectively\ }
\newcommand{\Subsection}{\subsection}
\providecommand{\nobreakdash}{\nobreak\hbox{-}}
\setlength{\emergencystretch}{2em}
\multlinegap0pt
\usepackage{bm}
\usepackage{bbm}
\usepackage{dsfont}
\usepackage{xspace}
\usepackage{cancel}
\usepackage{mathtools}
\usepackage{amsthm}
\makeatletter
\@ifundefined{theorem}{\newtheorem{theorem}{Theorem}}{}
\@ifundefined{lemma}{\newtheorem{lemma}[theorem]{Lemma}}{}
\makeatother
\makeatletter
\newcommand{\E}{\operatorname{\mathbb{E}}}
\newcommand{\Prob}{\operatorname{\Bbb{P}}}
\def\eps{\varepsilon}
\expandafter\ifx\csname proof\endcsname\relax
\newenvironment{proof}{\noindent{\bf Proof:}}{\qed}
\fi
\makeatother
\title[Stable matchings with switching costs]{Stable matchings with switching costs}
\author{Boris Pittel, Kirill Rudov}
\date{Source: arXiv:2507.14362v1 (CC BY 4.0)}
\begin{document}
\maketitle

\noindent\emph{Source and licence.} Stable matchings with switching costs. Version arXiv:2507.14362v1 (CC BY 4.0), distributed under the Creative Commons Attribution 4.0 International licence, as recorded in the arXiv licence metadata for this exact version and shown on the abstract page https://arxiv.org/abs/2507.14362v1. Full rights notice: licenses/arxiv-2507.14362-CC-BY-4.0.txt.\par\smallskip

\noindent\textit{Editorial scope.} Counted unit: the Lemma whose source label is \texttt{lem4} in the article (source0.tex lines 546--547, source label \texttt{lem4}), delivered together with the article's own complete original proof of it (source0.tex lines 549--684, one \texttt{proof} environment of 136 lines). No mathematics is written, repaired or re-derived: statement and proof are the article's own text line for line. Required context retained uncounted: A (lines 217--219); lem0 (lines 235--248); lem2 (lines 363--364). Every definition, display and label of those lines is retained unchanged. Omitted from this package and not redistributed: 1--216, 220--234, 249--362, 365--545, 548--548, 685--1137. Nothing inside a retained excerpt is omitted or altered: every retained body is delivered exactly as the article's lines are written, so no omission receipt of any kind accompanies this package. Citations: the retained excerpts carry no citation command at all, so no bibliography entry of the article is transcribed and no reference is redistributed. Reference adaptation: none was needed and none was made; every reference inside the delivered unit resolves inside it, so no unresolved reference or citation is produced. Printed slips kept literal: no printed word, symbol, hypothesis or number of the counted statement or of its proof was altered. Layout and class: the article's own class is \texttt{amsart} with packages including url, graphicx, bbold, latexsym, amsfonts, amsmath, amssymb, color; the wrapper is the lane's pinned \texttt{amsart} editorial wrapper at 10pt on A4, which restates the theorem-like environments the retained text needs and adds the article's own macro definitions that the retained text needs (\texttt{\textbackslash E}, \texttt{\textbackslash Prob}, \texttt{\textbackslash eps}), with extra wrapper packages (bm, bbm, dsfont, xspace, cancel, mathtools). The delivered numbering is the unit's own. Assets: no retained span contains a figure environment, a graphics-inclusion command, a PDF-page inclusion, a table or a TikZ picture; no image file is copied into this package, the package declares no asset dependency and no image file is redistributed with it. Licence: version arXiv:2507.14362v1 of this article is distributed under the Creative Commons Attribution 4.0 International licence, as recorded in the arXiv licence metadata for this exact version and shown on the abstract page https://arxiv.org/abs/2507.14362v1. A full rights notice is delivered at licenses/arxiv-2507.14362-CC-BY-4.0.txt. Characters: every retained excerpt is pure ASCII, so the delivered engine, XeLaTeX with its default Latin Modern text font, reports no missing character.
\begin{equation}\label{A}
 f_{\ell}(s,v_1,\dots,v_{\ell-1})=s^{\ell-1}\,\Bbb I\Bigl(\sum_{j<\ell}v_j\le 1\text{ and } \max_{j\le \ell}v_j\le\tfrac{1}{s}\Bigr),
\end{equation}

\begin{lemma}\label{lem0} As $\ell\to\infty$, in probability $\tfrac{L_{\ell}^+}{\ell^{-1}\log \ell}\to 1$, and $\ell\, U_{\ell}\to 2$. More sharply, for every $\rho>0$, we have 
\[
\Bbb P\bigl(L_{\ell}^+\le \ell^{-1}(\log\tfrac{\ell}{\log \ell}-\rho)\bigr)=O(\ell^{-d}), \quad
\forall\,d\in (0,e^{\rho}-1).,
\]
and for $\omega(\ell)\to\infty$ however slowly 
\[
\Bbb P\bigl(L_{\ell}^+\le \ell^{-1}\log(\ell\omega(\ell))\bigr)=(1+o(1))\exp\bigl(-\tfrac{1+o(1)}{\omega(\ell)}\bigr)\to 1, 
\]
while, for every $\delta<\tfrac{1}{3}$,
\[
\Bbb P\biggl(\big|\tfrac{\ell\, U_{\ell}}{2}-1\big|\ge \ell^{-\delta}\biggr)=O\bigl(e^{-\Theta(\ell^{1/3-\delta})}\bigr).
\]
\end{lemma}

\begin{lemma}\label{lem2} If $\eps \lambda =\omega(n^{-1}\log n)$, i.e. $\eps \lambda\gg n^{-1}\log n$, then $\E[S^{\eps}(n,n)]$ is super-polynomially large.
\end{lemma}

\begin{lemma}\label{lem4} If $\eps \lambda =\omega(n^{-1}\log n)$, i.e. $\eps \lambda\gg n^{-1}\log n$, then $\E[S^{\eps}(n,n+k)]$ is super-polynomially large, uniformly for all $k\ge 1$. %{\color{blue}(Remark: this NEW proof (in progress) seems to work for any $k\geq1$, not only $k=o(n)$.)
\end{lemma}

\begin{proof}
Since $\E[S^{\eps}(n,n+k)]$ is increasing in $\epsilon\lambda$, and $\epsilon\lambda=\omega(n^{-1}\log n)$, it is sufficient to focus only on $p=e^{-2\epsilon\lambda}=1-o(1)$ in our calculations below.

%{\color{magenta}
As before, we define the following variables:
\begin{equation*}
    s \equiv \sum_{i \in [n]}x_i,\quad s_j \equiv \sum_{i\neq j}x_i=s-x_j, \quad t \equiv \sum_{i \in [n]}x_i^2 , \quad \text{and} \quad t_j\equiv \sum_{i\neq j}x_i^2=t-x_j^2.
\end{equation*}
%}

Restrict attention to the set $D$ of all non-negative $x=(x_1,x_2,\ldots,x_n)$ such that
\begin{align}
\label{set_D_imbalance}
\begin{split}
&s_1(n)  \leq s \leq s_2(n), \\
&s^{-1} x_i \leq \frac{2\log n}{n},\quad i \in [n],\\
&s^{-2} t \leq \frac{3}{n},
\end{split}
\end{align}
where $s_2(n)>s_1(n)$ are to be defined later. (As before, the two bottom inequalities are essentially dictated by the connection between the fractions $s^{-1}x_i$ and the lengths $L_i$ of random subintervals 
of $[0,1]$, as well as Lemma \ref{lem0}.) According to the first two conditions, we have
\begin{equation*}
    x_i \leq s\frac{2\log n}{n} \leq \frac{2s_2(n)\log n}{n} \equiv \sigma_n.
\end{equation*}
In what follows, we choose $s_2(n)$ such that $\sigma_n \to 0$ (to be verified later). Thus, for sufficiently large $n$, we have $x_i\leq\sigma_n<1$ and $D \subset [0,1]^n$. Consequently, for such large $n$,
\begin{multline*}
P^\epsilon(n,n+k)\geq \tilde{P}^\epsilon(n,n+k) \equiv  \int_{\bold x \in D}\!\left(\prod_{j \in [n]}\left(\int_0^1 \prod_{i \neq j} (1-px_iy_j)\,dy_j\right)\right)\\
\times\prod_{\ell\in [n]}(1-\sqrt{p}x_{\ell})^k\,d\bold x.
\end{multline*}
Since $1-\alpha=\exp(-\alpha-\alpha^2(1+O(\alpha))/2)$ as $\alpha \to 0$, we have: for each $j$,
\begin{align*}
    \prod_{i \in[n], i\neq j} \left(1-px_iy_j\right) 
    &= \exp\left(-py_js_j-p^2y^2_jt_j\frac{1+O\left(\sigma_n\right)}{2}\right)\\
    &\geq \exp\left(-py_js-\frac{2p^2y^2_jt}{3}\right).
\end{align*}
By proceeding as in the proof of Lemma~\ref{lem2}, and using $t/s^2\leq3/n$, we bound
\begin{equation*}
\prod_{j \in [n]} \left( \int_0^1 \prod_{i \neq j} (1 - p x_i y_j) dy_j \right) \\
\geq e^{-4} \left(\frac{1 - e^{-ps}}{ps} \right)^n. 
\end{equation*}

Furthermore,
\begin{align*}
    \prod_{\ell\in [n]}(1-\sqrt{p}x_{\ell})^k 
    &= \exp\left(-k\sqrt{p}s-kpt\frac{1+O\left(\sigma_n\right)}{2}\right)\\
    &\geq \exp\left(-k\sqrt{p}s-\frac{2kps^2}{3}\frac{t}{s^2}\right)\\
    &\geq \exp\left(-k\sqrt{p}s-\frac{2kps^2}{n}\right)
\end{align*}
Based on the previous calculations, we have
\begin{align*}
P^\epsilon(n,n+k)\geq \int_{\bold x \in D}\!e^{-4} \left(\frac{1 - e^{-ps}}{ps} \right)^n \exp\left(-k\sqrt{p}s-\frac{2kps^2}{n}\right)\,d\bold x.
\end{align*}
We then switch to new variables $(u,v_1,v_2,\ldots,v_{n-1})$ defined as follows:
\begin{equation*}
    u \equiv \sum_{i \in [n]} x_i = s \qquad \text{and} \qquad v_i \equiv
    s^{-1}x_i, \quad i \in [n-1].
\end{equation*}
We also define $v_n \equiv s^{-1}x_n=1-\sum_{i \in [n-1]}v_i$. The conditions \eqref{set_D_imbalance} then become
\begin{align}
\begin{split}
&s_1(n)  \leq u \leq s_2(n), \\
&v_i \leq \frac{2\log(n)}{n},\quad i \in [n],\\
&\sum_{i \in [n]}v_i^2\leq \frac{3}{n},
\end{split}
\end{align}
Importantly, for these new variables, the resulting region is the direct product of the range of $u$ and the range of $\bold v$. By using the key equality \eqref{A},  we get then
\begin{align*}
\tilde{P}^\epsilon(n)& \geq \frac{e^{-4}p^{-n}}{(n-1)!} \left( \int_{s_1(n)}^{s_2(n)} (1 - e^{-pu})^n \exp\left(-k\sqrt{p}u\right)\cdot \frac{\exp\left(-\frac{2kpu^2}{n}\right)}{u} du \right)\\
&\quad\times \Prob\left(L^+_n  \leq \frac{2\log(n)}{n} \quad \text{and} \quad U_n \leq \frac{3}{n}\right),
\end{align*}
where $L^+_n \equiv  \max_{i \in [n]} L_i$ and $U_n \equiv \sum_{i \in [n]}L_i^2$. By Lemma \ref{lem0},  the probability of the event on the right-hand side tends to one as $n \to \infty$. Thus, it remains to examine
\begin{equation}\label{int*} 
\begin{aligned}
\int^\star& \equiv \int_{s_1(n)}^{s_2(n)}  \frac{\exp\left(H(u)-\frac{2kpu^2}{n}\right)}{u} du,\\
 H(u)& \equiv n\log (1 - e^{-pu})-k\sqrt{p}u.
\end{aligned}
\end{equation}
for appropriately chosen $s_{1,2}(n)$. Note that
\begin{equation*}
H'(u)=\frac{np}{e^{pu}-1}-k\sqrt{p},\,\, H''(u)=-\frac{np^2e^{pu}}{(e^{pu}-1)^2}
=-\frac{np^2}{(e^{pu/2}-e^{-pu/2})^2}\ge-\frac{n}{u^2},
\end{equation*}
i.e., the function $H(u)$ attains its maximum at 
$s(n) \equiv \frac{\log \left(\frac{n+k/\sqrt{p}}{k/\sqrt{p}}\right)}{p}$.
We anticipate that the dominant contribution to the integral $\int^\star$ is obtained when we set  
\begin{equation*}
 s_{1,2}(n)=s(n)(1\mp\delta(n))
\end{equation*}
for some $\delta(n) \to 0$. And for $u\in [s_1(n),s_2(n)]$, we have $H''(u)=-\Theta(n/s^2(n))$. (Note that $\sigma_n=\frac{2s_2(n)\log n}{n}\to 0$ for $n \to \infty$, as required.) Further, 
\begin{equation*}
H(s(n))=J(k/\sqrt{p}), \quad J(z):=n\log\frac{n}{n+z}+z\log\frac{z}{n+z},
\end{equation*}
where $J'(z)=\log\tfrac{z}{n+z}$, $J''(z)=\tfrac{n}{z(n+z)}>0$. It follows that 
\begin{equation*}
H(s(n))\ge n\log\frac{n}{n+k}+k\log\frac{k}{n+k}-\tfrac{1-p}{1+\sqrt{p}}k\log\frac{n+k}{k}.
\end{equation*}
So, expanding $H(u)$ in powers of $z:=u-s(n)$, we see that the integral $\int^*$ (defined in \eqref{int*}) is at least of order
\[
\frac{\exp(H(s(n))-3.1kps^2(n)/n)}{s(n)}\int_{-\delta(n)s(n)}^{\delta(n)s(n)}\exp\bigl(-z^2\Theta(n/s^2(n)\bigr)\,dz.\\
%=\tfrac{\exp(H(s(n))-3.1kps^2(n)/n)}{s(n)}\int_{-\infty}^{\infty}\exp\bigl(-z^2\Theta(n/s^2(n)\bigr)\,dz.
\]
Here the integral equals
\[
\tfrac{s(n)}{\sqrt{n}}\int_{-\delta(n)\sqrt{n}}^{\delta(n)\sqrt{n}}e^{-\Theta(\eta^2)}\,d\eta=\Theta(\tfrac{s(n)}{\sqrt{n}}),
\]
provided that $\delta (n)\sqrt{n}\to\infty$, which we assume. Therefore the integral $\int^*$ is at least of order
\begin{multline*}
\frac{\exp(H(s(n))-3.1kps^2(n)/n)}{\Theta(\sqrt{n})}\\
=\frac{n^n k^k}{(n+k)^{n+k}\Theta(\sqrt{n})}\exp\Bigl(-\tfrac{1-p}{1+\sqrt{p}}k\log\frac{n+k}{k}-3.1kps^2(n)/n\Bigr).
\end{multline*}
So, $\tilde{P}^\epsilon(n)$ is at least of order 
\begin{equation*}
\frac{p^{-n}}{(n-1)!}\frac{n^n k^k}{(n+k)^{n+k}\sqrt{n}}\exp\Bigl(-\frac{1-p}{1+\sqrt{p}}k\log\frac{n+k}{k}-3.1kps^2(n)/n\Bigr).
\end{equation*}
Combining this estimate with 
\begin{equation*}
\E[S^\epsilon (n,n+k)] 
\geq (n+k)_n \cdot \tilde{P}^\epsilon(n),\quad (n+k)_n\ge c e^{-k}(n+k)^{n+k} e^{-n},
%\binom{n+k}{n} =\Theta\bigl(\tfrac{(n+k)^{n+k}}{n^n k^k}\bigr),
\end{equation*}
$c$ being an absolute constant, we obtain that $\E[S^\epsilon (n,n+k)]$ is at least of order
\begin{equation*} 
%\tfrac{1}{\sqrt{n}}
\exp\Bigl(n\log(1/p)-\tfrac{1-p}{1+\sqrt{p}}k\log\tfrac{n+k}{k}-3.1kps^2(n)/n\Bigr).
\end{equation*}
Here, by the definition of $s(n)$,
\begin{equation*}
kps^2(n)/n=\frac{1}{\sqrt{p}}\cdot\tfrac{k}{n\sqrt{p}}\log^2\bigl(\frac{1+k/n\sqrt{p}}{k/n\sqrt{p}}\bigr)=O(1/\sqrt{p}),
\end{equation*}
since $\sup_{z>0} z\log^2\bigl(\frac{1+z}{z}\bigr)<\infty$.  In addition, $\frac{k}{n}\log\tfrac{n+k}{k}\le \frac{k}{n}\cdot \tfrac{n}{k}=1$, and $\log(1/p)>1-p$. Therefore $\E[S^\epsilon (n,n+k)]$ is at least of order
\begin{multline*}
\frac{1}{\sqrt{n}}\exp\Bigl(n\bigl(\log(1/p)-\frac{1-p}{1+\sqrt{p}}\bigr)+O(1/\sqrt{p})\Bigr)\\
\ge\frac{1}{\sqrt{n}} \exp\bigl(n(1-p)\frac{\sqrt{p}}{1+\sqrt{p}}+O(1/\sqrt{p})\bigr),
\end{multline*}
which is super-polynomially large for $p=e^{-2\eps\lambda}$, $\eps\lambda\to 0$, $\eps\lambda=\omega \tfrac{\log n}{n}$, $\omega=\omega(n)\to\infty$. This completes the proof of Lemma \ref{lem4}.
\end{proof}

\end{document}
